\section{Conclusion}
\label{sec:conclusion}

We believe the intersection of machine learning and economics presents a wide range of exciting research directions, and gives ample opportunity for new machine learning advances that will have significant positive social impact.
Our vision for the AI Economist is to enable an objective study of the impact of economic policies on real-world economies, at a level of complexity that traditional economics research cannot easily address. 

For tax policies in particular, we are hopeful that this kind of research can increase equality and productivity in the real world, helping to promote more just and healthy economies. We also hope that the AI Economist can foster transparency, reproducibility, and open and facts-based discussion about applying machine learning to economic decision-making, through our public research publications and open-source code. As such, we hope that future economic AI models can robustly and transparently augment real-world economic policy-making and in doing so improve social welfare.

In this paper, the economic agents and social planner were trained using model-free RL in AI-based, economic simulations.
A key benefit of using model-free RL is flexibility: for instance, any social objective can be used as the reward function for the planner. Moreover, it does not need any prior world knowledge to find a well-performing tax policy.
However, this approach assumes that the inputs and outputs to the agents' and planner's policy models are sufficient and well-defined. For instance, the planner should be able to observe all state information that is relevant for determining the optimal tax policy. Our initial experiments with human participants suggest that, in our problem setting, the state observed by the planner was sufficient to generalize well to human agents.
In future work, it would be interesting to explore which state information of the real world should be captured by economic simulators to enable generalization of policies from simulation to the real world.

This work is suggestive of the promise of AI-based, economic simulators for  learning economic policies with the potential to transfer well to the real world. 
We demonstrated that economic simulators can yield AI agents with economic behaviors that are consistent with economic intuition, for example agents that specialize as a consequence of their inherent skill level.
Our experiments also show that policies trained in such simulators can transfer well to settings with human participants, albeit for our limited problem setting.

Of course, these kinds of economic simulations still have many limitations. They do not yet model human-behavioral factors and interactions between people, including other-regarding utilities, and consider a relatively small economy.
Moreover, the concept of skill and the associated payoff, as used in our work, is still a limited representation of economic behavior in the real world. For instance, highly-skilled workers are not paid the same hourly wage across different industries, and skill might be hard to clearly define and measure clearly in certain professions.
Future simulations could improve the fidelity of simulated economic behavior by making use of real-world economic data, while advances in large-scale RL and engineering could increase the scope of economic simulations.


